\section{Tổng quan dự án}

\subsection{Mục tiêu}

Dự án này triển khai pipeline tạo ảnh hai giai đoạn dựa trên
\textit{Representation Autoencoders} (RAE). Ở thời điểm hiện tại, repo có hai
đường chạy chính:

\begin{itemize}[leftmargin=1.5em]
  \item \texttt{src/}: nhánh TorchXLA, tối ưu cho TPU với train và sampling
  Stage 2, reconstruction Stage 1, và FID phía host;
  \item \texttt{src\_jax/}: lớp tương thích JAX/NNX, giữ nguyên schema YAML của
  repo nhưng ánh xạ sang backend NNX để tránh nhân đôi toàn bộ codebase.
\end{itemize}

\subsection{Hai giai đoạn chính}

\subsubsection{Stage 1}

Stage 1 dùng một encoder biểu diễn đã được pretrain và đóng băng, kết hợp với
một decoder ViT để tái tạo ảnh. Đầu ra của Stage 1 là latent giàu ngữ nghĩa hơn
latent VAE thông thường.

\subsubsection{Stage 2}

Stage 2 học mô hình diffusion transformer trong không gian latent của Stage 1.
Thay vì học trực tiếp trên pixel, mô hình học quy luật biến đổi của latent, sau
đó giải mã latent sinh được thành ảnh RGB qua decoder của RAE.

\subsection{Phạm vi của branch hiện tại}

Nhánh hiện tại không chỉ có đường TorchXLA mà còn có adapter JAX mỏng. Phần được
hỗ trợ thực dụng nhất là:

\begin{itemize}[leftmargin=1.5em]
  \item huấn luyện Stage 2 bằng \texttt{src/train.py} trên TPU;
  \item huấn luyện Stage 2 bằng \texttt{src\_jax/train.py} trên JAX/NNX;
  \item sampling Stage 2 bằng cả \texttt{src/} và \texttt{src\_jax/};
  \item reconstruction Stage 1 để kiểm tra encoder-decoder trên cả hai đường;
  \item logging qua wandb;
  \item upload artifact hoặc workdir lên Hugging Face;
  \item FID offline, và online FID trong train loop ở đường XLA.
\end{itemize}

\subsection{Giới hạn hiện tại}

\begin{itemize}[leftmargin=1.5em]
  \item \texttt{src/} vẫn không có entrypoint riêng cho huấn luyện Stage 1 trên
  XLA;
  \item \texttt{src\_jax/} cố ý chưa port toàn bộ vòng huấn luyện Stage 1 kiểu
  GAN + LPIPS + discriminator, vì phần đó sẽ làm codebase lớn và khó bảo trì
  hơn đáng kể.
\end{itemize}

\subsection{Dòng dữ liệu cấp cao}

\begin{center}
\begin{tabular}{>{\raggedright\arraybackslash}p{0.24\linewidth} >{\raggedright\arraybackslash}p{0.64\linewidth}}
\toprule
Bước & Mô tả \\
\midrule
1 & Ảnh đầu vào được đưa qua encoder của Stage 1 \\
2 & Encoder sinh latent tensor có kích thước như \texttt{misc.latent\_size} \\
3 & Stage 2 học transport objective trên latent đó \\
4 & Khi sampling, latent được sinh bởi Stage 2 \\
5 & Decoder của Stage 1 giải mã latent thành ảnh RGB \\
\bottomrule
\end{tabular}
\end{center}

\subsection{Ý tưởng của đường JAX}

Đường JAX không tạo thêm một schema config mới. Thay vào đó:

\begin{itemize}[leftmargin=1.5em]
  \item vẫn đọc các block \texttt{stage\_1}, \texttt{stage\_2},
  \texttt{transport}, \texttt{sampler}, \texttt{guidance}, \texttt{misc},
  \texttt{training}, \texttt{eval};
  \item dùng \texttt{src\_jax/config\_adapter.py} để chuyển YAML hiện có sang
  config mà backend NNX hiểu được;
  \item bootstrap backend \texttt{diffuse\_nnx} theo commit cố định để branch
  gọn hơn, không phải chép toàn bộ framework vào repo.
\end{itemize}

\subsection{Tài liệu trong repo}

Ngoài PDF này, repo còn có bộ tài liệu markdown trong \texttt{docs/}. Phần
markdown thiên về runbook ngắn và tra cứu nhanh; PDF này gom lại kiến trúc,
workflow, cấu hình và vận hành cho cả đường XLA lẫn JAX.
